\scanpage{57}
polarisation donné et rigidification de Jacobi
\inserted{d'échelon} donné) en car $p$, est \deleted{\unsure{donc}}
une intersection complète, et de plus est \uline{loc.}
\uline{irréductible}, plus précisément
$\mathcal M_{\text{réd}}$ \uline{est \deleted{\unsure{intègre}}}
\inserted{lisse sur $k$}.
De plus, \deleted{la dimension de la \unsure{composante}
\unsure{correspondant} à une}

remontant en car. nulle et utilisant la platitude, on trouve que
\uline{la dimension du schéma modulaire} \inserted{\uline{envisagé}}
\uline{est bien $\dfrac{g(g+1)}2$ en chaque point}.
\margin{démonstration \unsure{indirecte}, on le voit
\unsure{directt} en car. $p$!}
Donc si la polarisation est de degré premier à $p$
\unsure{et donc} en car. quelconque, le schéma modulaire en car. $p$
est lisse.
(NB. il ne s'agit que de VA \uline{ordinaires})
\margin{67}
